\title{Group Sequence Policy Optimization}

\begin{abstract}
In this work, we present Group Sequence Policy Optimization (GSPO), a novel reinforcement learning approach designed for the stable, effective, and high-performance training of large language models. Distinct from prior methods that employ token-level importance ratios, GSPO utilizes importance ratios computed over entire sequences, incorporating sequence-level clipping, reward assignment, and policy updates. Our experiments show that GSPO surpasses the GRPO algorithm in both training efficiency and performance, provides pronounced stabilization for Mixture-of-Experts (MoE) reinforcement learning, and opens avenues for streamlining RL infrastructure design. These advantages have significantly advanced the performance of the latest Qwen3 models.
\end{abstract}

\section{Introduction}
\label{sec:intro}

Reinforcement learning (RL) has become a key approach for advancing the capabilities of language models at scale \citep{o1,dpsk-r1,qwq32b,qwen3}. Leveraging large-scale RL enables language models to address complex tasks, including competition-level mathematics and programming, by performing more intricate and extended reasoning.

For RL to be effectively scaled with increased computational resources, it is crucial to ensure that training remains both stable and resilient. Yet, cutting-edge RL algorithms such as GRPO \citep{grpo} frequently suffer from significant stability challenges when applied to training very large language models, which can lead to catastrophic and permanent failure of the model \citep{qwen3, minimax-m1}. These issues significantly impede further progress in enhancing language model abilities through ongoing RL training.

This study reveals that GRPO's instability arises from an incorrect application and subsequent breakdown of importance sampling weights within its algorithm. This flaw introduces substantial training noise characterized by high variance, which escalates with longer response sequences and is further intensified by the clipping procedure, eventually leading to model collapse.

To overcome these fundamental shortcomings, we introduce \textbf{Group Sequence Policy Optimization (GSPO)}, a novel reinforcement learning algorithm designed for the efficient training of large language models. GSPO's main contribution is a rigorously justified definition of the importance ratio, derived from sequence likelihood \citep{click}, which closely adheres to the underlying principle of importance sampling. Moreover, GSPO calculates normalized rewards as the advantages across several generated responses per query, facilitating a robust connection between sequence-level reward signals and policy optimization.

Through our empirical analysis, we find that GSPO markedly outperforms GRPO in terms of training stability, efficiency, and overall effectiveness. Importantly, GSPO intrinsically addresses the instability issues commonly encountered during reinforcement learning with sizable Mixture-of-Experts (MoE) models, thereby obviating the necessity for intricate stabilization procedures and offering prospects for streamlining RL system architecture. These advantages of GSPO have played a pivotal role in the remarkable gains observed in the recent Qwen3 models. Looking ahead, we anticipate that GSPO will serve as a reliable and scalable algorithmic base for advancing the large-scale RL training landscape in language models.

\section{Preliminaries}

\paragraph{Notation}
Throughout this work, we refer to an autoregressive language model with parameters $\theta$ as the policy $\pi_\theta$. The variable $x$ represents a query, while $\mathcal{D}$ denotes the collection of queries. For any query $x$, a corresponding response $y$ is evaluated for its probability according to the policy $\pi_\theta$, specified as $\pi_\theta (y | x)=\prod_{t=1}^{|y|} \pi_\theta (y_t | x, y_{<t} )$, where $|y|$ counts the tokens in $y$. Each pair $(x, y)$ is assessed by a verifier $r$, yielding a reward $r(x, y)$ within the interval $[0, 1]$.

\paragraph{Proximal Policy Optimization (PPO)} Using samples generated from the old policy $\pi_{\theta_\text{old}}$, PPO \citep{ppo} constrains the policy update within a proximal region of the old policy through the clipping mechanism. Specifically, PPO employs the following objective for policy optimization (we omit the KL regularization term hereinafter for brevity, as it is not the focus of this paper):

\begin{align}
\mathcal{J}_\text{PPO}(\theta) = \mathbb{E}_{ x \sim \mathcal{D},\, y \sim \pi_{\theta_\text{old}}( \cdot | x) }
\left[ \frac{1}{|y|} \sum_{t=1}^{|y|} 
\min \left( w_{t}(\theta) \widehat{A}_{t},  \, \mathrm{clip} \left( w_{t}(\theta), 1 - {\varepsilon}, 1 + {\varepsilon}\right) \widehat{A}_{t} \right)
\right],
\end{align}

Here, the importance ratio for the token $y_t$ is given as $
w_{t}(\theta) = \frac{ \pi_{\theta} (y_{t} | x, y_{<t}) }{ \pi_{\theta_\text{old}} (y_{t} | x,y_{<t})} $. The advantage $\widehat{A}_{t}$ associated with $y_t$ is computed by an auxiliary value model, and $\varepsilon$ denotes the threshold for clipping the importance ratios.

A principal obstacle faced by PPO in practical applications is its significant dependence on the value model. The value model is typically architected to be as large as the policy model, resulting in substantial demands on both memory and compute resources. Additionally, the algorithm's performance is closely tied to the precision of the value estimates produced. Developing an accurate value model is difficult in itself, but extending its accuracy to longer output sequences or more intricate tasks becomes increasingly problematic.

\paragraph{Group Relative Policy Optimization (GRPO)} GRPO \citep{grpo} bypasses the need for the value model by computing the relative advantage of each response within a group of responses to the same query. Specifically, GRPO optimizes the following objective:

\begin{align}
\label{equ:grpo}
\mathcal{J}_\text{GRPO}(\theta) = \mathbb{E}_{ x \sim \mathcal{D},\, \{y_i\}_{i=1}^G \sim \pi_{\theta_\text{old}}( \cdot | x) }
\left[ \frac{1}{G} \sum_{i=1}^{G} \frac{1}{|y_i|} \sum_{t=1}^{|y_i|} 
\min \left( w_{i,t}(\theta) \widehat{A}_{i,t},  \, \mathrm{clip} \left( w_{i,t}(\theta), 1 - {\varepsilon}, 1 + {\varepsilon}\right) \widehat{A}_{i,t} \right)
\right],
\end{align}

where $G$ is the number of generated responses to each query $x$ (i.e., the group size), and the importance ratio $w_{i,t}(\theta)$ and advantage $\widehat{A}_{i,t}$ of token $y_{i,t}$ are:

\begin{align}
    w_{i,t}(\theta)=\frac{ \pi_{\theta} (y_{i,t} | x, y_{i,<t}) }{ \pi_{\theta_\text{old}} (y_{i,t} | x,y_{i,<t})},\quad\ 
    \widehat{A}_{i,t} = \widehat{A}_{i} = \frac{r(x, y_i) - \mathrm{mean} \left( \{ r(x, y_i) \}_{i=1}^G \right) }{ \mathrm{std} \left( \{ r(x, y_i) \}_{i=1}^G \right) },
\end{align}

respectively, where all the tokens in $y_i$ share the same advantage as $\widehat{A}_{i}$.

\section{Motivation}
\label{sec:motivation}

As models become larger and more sparse (for instance, in the case of Mixture-of-Experts architectures) and as response sequences lengthen, achieving efficient hardware usage in RL training often requires substantial rollout batch sizes. To enhance sample efficiency, it is customary to divide these large rollout batches into several smaller mini-batches for performing gradient updates. However, this approach inherently gives rise to an off-policy learning scenario, since responses $y$ are generated under a prior policy $\pi_{\theta_\text{old}}$ instead of the current policy $\pi_\theta$ targeted by optimization. Consequently, mechanisms such as clipping in PPO and GRPO are vital, as they limit the influence of data that are too far off-policy in the computation of policy gradients.

While mechanisms like clipping aim to manage this off-policy discrepancy, we identify a more fundamental issue in GRPO: \textit{its objective is ill-posed}. This problem becomes particularly acute when training large models on long-response tasks, leading to catastrophic model collapse. The ill-posed nature of the GRPO objective stems from a misapplication of importance sampling weights. The principle of importance sampling is to estimate the expectation of a function $f$ under a target distribution $\pi_\text{tar}$ by re-weighting samples drawn from a behavior distribution $\pi_\text{beh}$:

\begin{align}
\mathbb{E}_{z \sim \pi_\text{tar}} \left[ f(z) \right] = \mathbb{E}_{z \sim \pi_\text{beh}} \left[ \frac{ \pi_\text{tar} (z) }{ \pi_\text{beh} (z)} \, f(z) \right].
\end{align}

A key aspect here is the requirement to take an average over a large number of samples ($N \gg 1$) drawn from the behavior policy $\pi_\text{beh}$, so that the importance weight $\frac{ \pi_\text{tar} (z) }{ \pi_\text{beh} (z)}$ properly compensates for the discrepancy between the distributions.

On the other hand, GRPO calculates the importance weight $\frac{ \pi_{\theta} (y_{i,t} | x, y_{i,<t}) }{ \pi_{\theta_\text{old}} (y_{i,t} | x,y_{i,<t})}$ at each token step $t$. This approach uses just a single token $y_{i,t}$ sampled from the next-token distribution $\pi_{\theta_\text{old}}( \cdot | x, y_{i, <t})$ per position, failing to enact the desired correction between distributions. Rather than reducing mismatch, it introduces substantial variance in the gradient estimates, which becomes more pronounced throughout extended sequences and is further intensified by the clipping procedure. Empirical results indicate that such variance can precipitate catastrophic model collapse, typically unrecoverable. Once collapse ensues, attempts to resume training—including restoring prior checkpoints, extensive hyperparameter (e.g., clipping range) adjustments, increasing generated sequence length, or altering RL queries—prove ineffective.

These findings illuminate a fundamental flaw in GRPO's methodology. The inefficacy of the token-level importance weighting underscores a central insight: \textbf{the granularity of the optimization objective must align with the reward's granularity}. Because rewards are assigned at the full-sequence level, applying off-policy corrections at individual tokens presents inherent challenges. This insight compels us to abandon the token-level framework and consider applying importance weighting and optimization directly at the \textit{sequence level}.

\section{Algorithm}

\subsection{GSPO: Group Sequence Policy Optimization}

Although the token-level importance weight $\frac{ \pi_{\theta} (y_{i,t} | x, y_{i,<t}) }{ \pi_{\theta_\text{old}} (y_{i,t} | x,y_{i,<t})}$ introduces difficulties in GRPO, we note that, for language generation tasks, the \textit{sequence-level} importance weight $\frac{ \pi_{\theta} (y | x) }{ \pi_{\theta_\text{old}} (y | x)}$ possesses a well-defined theoretical interpretation. Specifically, it quantifies the extent to which a response $y$—drawn from $\pi_{\theta_\text{old}} (\cdot | x)$—differs from $\pi_{\theta} (\cdot | x)$, thereby directly corresponding to sequence-level reward calculations and providing an interpretable basis for implementing a clipping mechanism.

Based on this straightforward observation, we propose the \textbf{Group Sequence Policy Optimization (GSPO)} algorithm. GSPO employs the following sequence-level optimization objective:

\begin{align}
\mathcal{J}_\text{GSPO} (\theta) =
\mathbb{E}_{ x \sim \mathcal{D},\, \{y_i\}_{i=1}^G \sim \pi_{\theta_\text{old}}( \cdot | x) }
\left[ 
\frac{1}{G} \sum_{i=1}^{G}
\min \left( s_{i}(\theta)  \widehat{A}_{i},  \, \mathrm{clip} \left( s_{i}(\theta), 1 - {\varepsilon}, 1 + {\varepsilon} \right) \widehat{A}_{i} \right) 
\right],
\label{equ:gspo}
\end{align}

where we adopt the group-based advantage estimation:

\begin{align}
\widehat{A}_{i} = \frac{r(x, y_i) - \mathrm{mean} \left( \{ r(x, y_i) \}_{i=1}^G \right) }{ \mathrm{std} \left( \{ r(x, y_i) \}_{i=1}^G \right) },
\end{align}

and define the importance ratio $s_{i}(\theta)$ based on sequence likelihood \citep{click}:

\begin{align}
s_{i}(\theta) = \left( \frac{ \pi_{\theta} (y_i | x) }{ \pi_{\theta_\text{old}} (y_i | x)} \right)^{\frac{1}{|y_i|}}
=
\exp \left( \frac{1}{|y_i|} \sum_{t=1}^{|y_i|} \log \frac{ \pi_{\theta} (y_{i,t} | x, y_{i,<t}) }{ \pi_{\theta_\text{old}} (y_{i,t} | x,y_{i,<t})} \right).
\end{align}

Consequently, GSPO employs response-level clipping, rather than token-level clipping, in order to remove excessively ``off-policy'' samples from gradient calculation. This strategy aligns with both reward assignment and optimization at the sequence level. We implement length normalization in $s_{i}(\theta)$, which mitigates variance and confines $s_{i}(\theta)$ to a consistent numerical interval. Without normalization, changes in the probability of a handful of tokens could lead to substantial volatility in the sequence-level importance ratio. Additionally, importance ratios for responses of varying lengths would necessitate different clipping thresholds. Furthermore, it should be observed that the clipping boundaries in GSPO, compared to earlier methods such as GRPO, typically differ by orders of magnitude, attributable to differences in how importance ratios are defined.

\subsection{Gradient Analysis}
\label{subsec:gradient}

We can derive the gradient of the GSPO objective as follows (clipping is omitted for brevity):

\begin{align}
\nabla_{\theta} \mathcal{J}_\text{GSPO} (\theta)
=&\ 
\nabla_{\theta} \mathbb{E}_{ x \sim \mathcal{D},\, \{y_i\}_{i=1}^G \sim \pi_{\theta_\text{old}}( \cdot | x) }
\left[ 
\frac{1}{G} \sum_{i=1}^{G} s_{i}(\theta) \widehat{A}_{i}
\right] \\
= &\ 
\mathbb{E}_{ x \sim \mathcal{D},\, \{y_i\}_{i=1}^G \sim \pi_{\theta_\text{old}}( \cdot | x) }
\left[ 
\frac{1}{G} \sum_{i=1}^{G}
s_{i}(\theta) \widehat{A}_{i}
\cdot \nabla_{\theta} \log s_{i}(\theta)
\right] \\
=&\ 
\mathbb{E}_{ x \sim \mathcal{D},\, \{y_i\}_{i=1}^G \sim \pi_{\theta_\text{old}}( \cdot | x) }
\left[ 
\frac{1}{G} \sum_{i=1}^{G} 
\left( \frac{ \pi_{\theta} (y_i | x) }{ \pi_{\theta_\text{old}} (y_i | x)} \right)^{\frac{1}{|y_i|}} \widehat{A}_{i} 
\cdot \frac{1}{|y_i|} \sum_{t=1}^{|y_i|} \nabla_{\theta} \log \pi_{\theta} (y_{i,t} | x, y_{i,<t}) 
\right].
\label{equ:gspo_grad}
\end{align}

For comparison, the gradient of the GRPO objective is as follows (note that $\widehat{A}_{i,t} = \widehat{A}_{i}$):

\begin{align}
\nabla_{\theta} \mathcal{J}_\text{GRPO} (\theta) 
=&\ 
\nabla_{\theta} \mathbb{E}_{ x \sim \mathcal{D},\, \{y_i\}_{i=1}^G \sim \pi_{\theta_\text{old}}( \cdot | x) }
\left[ \frac{1}{G} \sum_{i=1}^{G} \frac{1}{|y_i|} \sum_{t=1}^{|y_i|} 
w_{i,t}(\theta) \widehat{A}_{i,t}
\right] \\
=&\
\mathbb{E}_{ x \sim \mathcal{D},\, \{y_i\}_{i=1}^G \sim \pi_{\theta_\text{old}}( \cdot | x) }
\left[ \frac{1}{G} \sum_{i=1}^{G} \widehat{A}_{i} 
\cdot \frac{1}{|y_i|} \sum_{t=1}^{|y_i|} 
\frac{ \pi_{\theta} (y_{i,t} | x, y_{i,<t}) }{ \pi_{\theta_\text{old}} (y_{i,t} | x,y_{i,<t})} 
\nabla_{\theta} \log \pi_{\theta} (y_{i,t} | x, y_{i,<t})  
\right].
\end{align}

Thus, the principal difference between GSPO and GRPO is rooted in \textit{the method each approach employs to assign weights to the gradients derived from the token log likelihoods}. GRPO utilizes per-token ``importance weights'' given by $\frac{ \pi_{\theta} (y_{i,t} | x, y_{i,<t}) }{ \pi_{\theta_\text{old}} (y_{i,t} | x,y_{i,<t})}$, resulting in non-uniform weighting. These weights can differ significantly, falling within the range $(0, 1+\varepsilon]$ for cases where $\widehat{A}_{i} >0$, or $[1-\varepsilon, +\infty)$ when $\widehat{A}_{i} < 0$, making their influence potentially substantial and cumulative, possibly causing unpredictable effects as training proceeds. By contrast, GSPO applies uniform weighting across all tokens in a given response, thereby removing the instability associated with the varying weights in GRPO.

\subsection{GSPO-token: A Token-level Objective Variant}

In scenarios like multi-turn RL, we may desire a finer-grained advantage adjustment than the sequence level. To this end, we introduce a token-level objective variant of GSPO, namely \textbf{GSPO-token}, to allow token-wise advantage customization:

\begin{align}
\mathcal{J}_\text{GSPO-token}(\theta)
=
\mathbb{E}_{ x \sim \mathcal{D},\, \{y_i\}_{i=1}^G \sim \pi_{\theta_\text{old}}( \cdot | x) }
\left[ 
\frac{1}{G} \sum_{i=1}^{G} \frac{1}{|y_i|} \sum_{t=1}^{|y_i|} 
\min \left( s_{i,t}(\theta) \widehat{A}_{i,t},  \, \mathrm{clip} \left( s_{i,t}(\theta), 1 - {\varepsilon}, 1 + {\varepsilon} \right) \widehat{A}_{i,t} \right) 
\right],
\label{equ:gspo-token}
\end{align}

where

\begin{align}
s_{i,t}(\theta) 
= \mathrm{sg} \left[ s_{i}(\theta) \right]  \cdot \frac{ \pi_{\theta} (y_{i,t} | x, y_{i,<t}) }{ \mathrm{sg} \left[ \pi_{\theta} (y_{i,t} | x, y_{i,<t}) \right] },
\end{align}

and $\mathrm{sg}[\cdot]$ denotes only taking the numerical value but stopping the gradient, corresponding to the \texttt{detach} operation in PyTorch. The gradient of GSPO-token can be derived as:

\begin{align}
\nabla_{\theta} \mathcal{J}_\text{GSPO-token}(\theta)
=&\ 
\nabla_{\theta} \mathbb{E}_{ x \sim \mathcal{D},\, \{y_i\}_{i=1}^G \sim \pi_{\theta_\text{old}}( \cdot | x) }
\left[ 
\frac{1}{G} \sum_{i=1}^{G}
\frac{1}{|y_i|} \sum_{t=1}^{|y_i|} 
s_{i,t}(\theta) \widehat{A}_{i,t}
\right] \\
=&\ 
\mathbb{E}_{ x \sim \mathcal{D},\, \{y_i\}_{i=1}^G \sim \pi_{\theta_\text{old}}( \cdot | x) }
\left[ 
\frac{1}{G} \sum_{i=1}^{G} s_i (\theta) \cdot \frac{1}{|y_i|} \sum_{t=1}^{|y_i|} 
\widehat{A}_{i,t} \frac{ \nabla_{\theta} \pi_{\theta} (y_{i,t} | x, y_{i,<t}) }{ \pi_{\theta} (y_{i,t} | x, y_{i,<t}) }
\right] \\
=&\ 
\mathbb{E}_{ x \sim \mathcal{D},\, \{y_i\}_{i=1}^G \sim \pi_{\theta_\text{old}}( \cdot | x) }
\left[ 
\frac{1}{G} \sum_{i=1}^{G}  \left( \frac{ \pi_{\theta} (y_i | x) }{ \pi_{\theta_\text{old}} (y_i | x)} \right)^{\frac{1}{|y_i|}}
\cdot \frac{1}{|y_i|} \sum_{t=1}^{|y_i|}
\widehat{A}_{i,t} \nabla_{\theta} \log \pi_{\theta} (y_{i,t} | x, y_{i,<t})  
\right].
\label{equ:gspo-token_grad}
\end{align}

Observe that the expression $\frac{ \pi_{\theta} (y_{i,t} | x, y_{i,<t}) }{ \mathrm{sg} \left[ \pi_{\theta} (y_{i,t} | x, y_{i,<t}) \right] }$ evaluates to 1, which implies that $s_{i,t}(\theta)$ takes the same numerical value as $s_{i}(\theta)$. When we analyze Equation~\eqref{equ:gspo} alongside Equation~\eqref{equ:gspo-token}, as well as Equation~\eqref{equ:gspo_grad} in comparison to Equation~\eqref{equ:gspo-token_grad}, it becomes evident that GSPO-token and GSPO are mathematically equivalent in terms of their optimization objective, clipping rule, and theoretical gradient, provided that the advantage estimates for all tokens in $y_i$ are held constant (specifically, $\widehat{A}_{i,t} = \widehat{A}_{i}$). Nevertheless, GSPO-token retains additional adaptability, as it permits the advantage values to be tailored individually for each token.

\section{Experiments and Discussion}

\subsection{Empirical Results}

We conduct experiments utilizing a cold-start model derived from fine-tuning Qwen3-30B-A3B-Base, presenting both training reward trajectories and evaluation results across the AIME'24 (mean Pass@1 using 32 samplings), LiveCodeBench (202410-202502, mean Pass@1 over 8 samplings), and CodeForces (Elo Rating) test suites. During reinforcement learning, rollout batches are divided into four mini-batches for separate gradient updates. For GSPO, we configure the left and right clipping thresholds in Equation~\eqref{equ:gspo} as 3e-4 and 4e-4, respectively. For comparative purposes, we employ GRPO as the baseline approach and tune its left and right clipping values in Equation~\eqref{equ:grpo} to 0.2 and 0.27, respectively, to facilitate an equitable evaluation. It is important to highlight that GRPO requires the use of the Routing Replay strategy to ensure proper convergence of MoE RL, as further addressed in \S~\ref{subsec:routing_replay}. In contrast, \textit{GSPO dispenses with this necessity}.

Figure~\ref{fig:results} illustrates that GSPO enables consistent and stable progress during training. Notably, \textbf{GSPO facilitates sustained enhancement in model performance by increasing training compute, frequently refreshing the query set, and lengthening the generation outputs}. Additionally, GSPO outperforms GRPO in terms of training efficiency, yielding higher accuracy and benchmark results given equivalent training compute and query consumption. Lastly, our application of GSPO to RL training on recent Qwen3 models confirms its effectiveness, decisively demonstrating the ability of GSPO to harness RL scaling in large language models.

\subsection{Curious Observation on Clipping Fractions}

One notable difference between GSPO and GRPO lies in the approach to clipping: GSPO applies clipping at the response level, while GRPO operates at the token level. As illustrated in Figure~\ref{fig:clipping}, GSPO clips a substantially greater proportion of tokens—by about two orders of magnitude—compared to GRPO, and altering the clipping thresholds does not diminish this marked discrepancy. Surprisingly, despite GSPO excluding more tokens from the training (or gradient calculation) process, it attains greater training efficiency relative to GRPO. This seemingly paradoxical result—where the removal of a higher proportion of tokens corresponds with improved training outcomes—underscores the fact that GRPO's token-wise gradient estimations tend to be intrinsically noisy and inefficient for leveraging sampled data. Conversely, GSPO's strategy of sequence-level clipping enables it to generate a more robust and informative learning signal.

\subsection{Benefit of GSPO for MoE Training}
\label{subsec:routing_replay}

\paragraph{Background}
In contrast to RL training with dense models, the inherently sparse activations of MoE architectures present distinct challenges regarding stability. Specifically, we observe that when utilizing the GRPO algorithm, MoE models can exhibit pronounced \textit{expert-activation volatility}, which hinders the convergence of RL training processes. More concretely, following one or several rounds of gradient-based updates, the subset of experts selected for generating the same response may vary considerably. As an illustration, with the 48-layer Qwen3-30B-A3B-Base model, every RL gradient update causes approximately 10\% of the experts chosen for the same rollout sample under the updated policy $\pi_{\theta}$ to differ from those selected with the previous policy $\pi_{\theta_\text{old}}$. This effect, especially marked in deeper MoE architectures, induces substantial fluctuations in the token-level importance ratios $w_{i,t}(\theta) = \frac{ \pi_{\theta} (y_{i,t} | x, y_{i,<t}) }{ \pi_{\theta_\text{old}} (y_{i,t} | x,y_{i,<t})}$ and undermines their validity, as elaborated in \S~\ref{sec:motivation} and~\ref{subsec:gradient}. Such instability impedes the proper convergence of RL optimization.

\paragraph{Our Previous Approach} In our earlier work, we utilized the \textbf{Routing Replay} training method to address this issue. Our approach involved storing the experts activated by $\pi_{\theta_\text{old}}$, then subsequently using these same routing paths when evaluating $\pi_{\theta}$ while calculating the importance ratios $w_{i,t}(\theta) = \frac{ \pi_{\theta} (y_{i,t} | x, y_{i,<t}) }{ \pi_{\theta_\text{old}} (y_{i,t} | x,y_{i,<t})}$. This process ensures that, for every token $y_{i,t}$, both $\pi_{\theta} (y_{i,t} | x, y_{i,<t})$ and $\pi_{\theta_\text{old}} (y_{i,t} | x,y_{i,<t})$ operate on identical expert networks, thus preserving the consistency and stability of importance ratios at the token level, and maintaining optimization over a consistent set of activated networks during gradient updates. Figure~\ref{fig:routing_replay} illustrates how Routing Replay is crucial for achieving stable convergence in GRPO training for MoE models.

\paragraph{Benefit of GSPO} While Routing Replay facilitates the convergence of GRPO training in MoE models, its reliance on reusing routing modes imposes extra demands on memory and communication, and it may restrict the model's effective capacity. Conversely, as illustrated in Figure~\ref{fig:results}, GSPO removes the necessity for Routing Replay, computing importance ratios $s_i(\theta)$ in the standard manner. GSPO achieves stable optimization and normal convergence without additional complexity. The central insight is that GSPO prioritizes the likelihood of the sequence, specifically $\pi_{\theta} (y_i | x)$, and does not depend on the likelihood of individual tokens, $\pi_{\theta} (y_{i,t} | x, y_{i,<t})$. Because the MoE model retains its core language modeling abilities, the sequence-level likelihood remains consistent and avoids sudden changes. In effect, GSPO directly addresses the volatility in expert activation seen in MoE models and eliminates the need for intricate solutions such as Routing Replay. This leads to a streamlined and robust training procedure, empowering the model to fully exploit its inherent capacity.

\subsection{Benefit of GSPO for RL Infrastructure}

Due to differences in numerical precision between training frameworks such as Megatron and inference frameworks like SGLang and vLLM, it is common practice to recompute the likelihoods of sampled outputs under the prior policy $\pi_{\theta_\text{old}}$ using the training engine. Nevertheless, because GSPO relies on sequence-level likelihoods—rather than those calculated at each token—it is inherently less sensitive to such precision mismatches. As a result, GSPO permits direct utilization of likelihood estimates from the inference engine during optimization, eliminating the overhead of recomputation through the training engine. This feature is particularly advantageous in applications involving partial rollout, multi-turn RL, and systems that separate training from inference.

\section{Conclusion}

Group Sequence Policy Optimization (GSPO) is introduced as a novel reinforcement learning approach aimed at enhancing large language model training. GSPO leverages the principle of importance sampling, establishing importance ratios derived from sequence likelihood, and implements sequence-level clipping, rewarding, and optimization techniques. Compared to GRPO, GSPO achieves substantially greater stability, efficiency, and performance during training, and proves especially effective for conducting large-scale RL training of MoE models—serving as a critical factor in the major advancements seen in the most recent Qwen3 models. Positioned as a scalable framework, GSPO will support further expansion of RL methods and is expected to drive key progress in artificial intelligence.

\end{document}